% Paper 6 (GRPO-Registry) query API + auditor badge section
\section{Query API, Auditor Badge, and \texttt{stackdiff}}
\label{sec:p6-query}

The reference CLI (\texttt{registry/query.py}, standard library only) is a
specification by example: any conforming implementation must reproduce its
outputs on the shipped entries.

\subsection{Query API}
\label{sec:p6-query-api}
The query surface is deliberately small---four verbs over the entry
directory:

\begin{lstlisting}[basicstyle=\ttfamily\scriptsize, frame=single,
  caption={The reference CLI surface.}, label={lst:p6-cli}]
python3 registry/query.py list
python3 registry/query.py query --item reference_kl [--full]
python3 registry/query.py badge [entry_id]
python3 registry/query.py stackdiff <entry_a> <entry_b>
\end{lstlisting}

\texttt{query --item} classifies every stack as \texttt{FULL},
\texttt{PARTIAL}, or \texttt{UNREPORTED} on one MIN-REPORT item, where the
classification counts non-null leaves; the null-vs-explicit-absence convention
of Section~\ref{sec:p6-schema} does the real work here. \texttt{list} and the
JSON records themselves are the escape hatch: the format is plain JSON, so
\texttt{jq}, pandas, or a five-line script can express any richer query.

\subsection{The MIN-REPORT Auditor Badge}
\label{sec:p6-badge}
The badge maps a stack record to $[0,100]$:
\begin{equation}
  \mathrm{badge}(s) \;=\;
  \Big\lfloor 100 \cdot \tfrac{1}{7}\textstyle\sum_{i=1}^{7}
  c_i(s) \Big\rceil,
  \qquad
  c_i(s) = \frac{\#\{\text{non-null leaves of item } i\}}
                {\#\{\text{leaves of item } i\}},
  \label{eq:p6-badge}
\end{equation}
i.e.\ the mean per-item reporting coverage. Two design decisions matter.
First, items are weighted equally: a stack cannot buy a high badge by
over-documenting its sampler while staying silent on its loss form. Second,
the badge is \emph{outcome-blind}: the \texttt{outcomes} block is excluded, so
a collapsed run and a converged run with equal reporting earn equal badges.
The badge answers ``can this stack be audited?''---never ``is this stack
good?''. On the seed entries the badge spans 43 (the zvf-iter130
risk-panel entries, which omit loss-form, KL, and group-size-schedule
disclosure) to 96 (open trainer with explicit negative declarations),
with the managed-runtime head-to-head at 57 and the open-framework
config dumps at 60--74 in between (Table~\ref{tab:p6-entries}).

\subsection{\texttt{stackdiff} and the R0--R5 Flip-Risk Ladder}
\label{sec:p6-stackdiff}
\texttt{stackdiff} compares two stack records and reports every differing
leaf, each tagged with a rung on a flip-risk ladder; the verdict is the
highest rung triggered:

\begin{itemize}[leftmargin=1.5em, itemsep=0.1em]
  \item \textbf{R0} --- identical manifests.
  \item \textbf{R1} --- nuisance-only differences (metadata, logging).
  \item \textbf{R2} --- runtime differences: backend, precision, reference
        placement. Not benign---our same-label audit contains R2 differences
        plus a base-checkpoint mismatch---but auditable when fully reported.
  \item \textbf{R3} --- loss-form or KL-handling differences: the update rule
        itself differs.
  \item \textbf{R4} --- the variant-delta sets differ: a technique credited to
        the shared label is \texttt{absent} or merely \texttt{surrogate} on
        one side. Any cross-stack comparison under the shared label is now at
        risk of a label flip.
  \item \textbf{R5} --- unauditable: a differing item is null (unreported) on
        at least one side, so the flip risk cannot even be bounded.
\end{itemize}

The ladder is intentionally a \emph{risk} scale, not an effect-size scale: our
largest observed same-label swing ($17\times$) contained R2 differences and a
checkpoint mismatch, so it cannot estimate an R2 effect. Such differences are
diagnosable from complete manifests; R5 says the manifests themselves are
insufficient. When the verdict reaches R4 or above, the CLI prints an explicit
warning that outcome differences must not be attributed to the shared
algorithm label. A companion spec-level tool family (a trainer plugin that
auto-emits the seven-field run manifest during training, and a LeverTrace check
that asks whether a cited paper actually reports the lever it is credited for)
shares this ladder; only \texttt{stackdiff} and the badge are implemented in
the shipped CLI, and we scope claims accordingly.
